\documentclass[11pt,a4paper]{article}
\usepackage[T1]{fontenc}
\usepackage[utf8]{inputenc}
\usepackage[spanish,es-nodecimaldot]{babel}
\usepackage{lmodern}
\usepackage{geometry}
\usepackage{microtype}
\usepackage{graphicx}
\usepackage{booktabs}
\usepackage{multirow}
\usepackage{amsmath}
\usepackage{upquote}
\usepackage{xcolor}
\usepackage{pgfplots}
\usepackage[hidelinks]{hyperref}
\pgfplotsset{compat=1.18}
\geometry{top=2.3cm,bottom=2.5cm,left=2.5cm,right=2.5cm}
\setlength{\parskip}{0.35em}
\setlength{\parindent}{0pt}
\newcommand{\ruta}[1]{\path{#1}}
\definecolor{serieroaring}{HTML}{2A78D6}
\definecolor{seriescan}{HTML}{EB6834}
\definecolor{seriegris}{HTML}{9A9A94}
\begin{document}
\begin{titlepage}
\centering
\includegraphics[width=3.5cm]{logo_utec.png/images.png}\par
\vspace{1cm}
Universidad de Ingeniería y Tecnología\\
CS3014 Estructuras de Datos Avanzadas\par
\vspace{1.8cm}
{\LARGE\bfseries Roaring Bitmaps para un índice invertido de registros de ventas: implementación y evaluación\par}
\vspace{1cm}
{\large Informe final\par}
\vspace{1.5cm}
\textbf{Integrantes}\\[0.3cm]
David Herencia Galván\\
Sofía Herrera Salazar\par
\vfill
\end{titlepage}

\begin{abstract}
Se implementó en C++17 una versión propia de Roaring Bitmaps para conjuntos de enteros de 32 bits y se aplicó como índice invertido sobre las 541\,909 filas del conjunto Online Retail. La estructura divide cada entero en una parte alta de 16 bits, que selecciona un contenedor, y una parte baja de 16 bits, que se guarda en un arreglo ordenado (hasta 4096 valores) o en un bitmap de 8192 bytes (más de 4096 valores). Sobre esta representación se implementaron inserción, pertenencia, intersección (AND) y unión (OR), con un algoritmo específico para cada par de tipos de contenedor. El índice contiene 2310 términos y 38 países; sus postings ocupan 5,76 MB en el modelo de espacio usado, el 49\% de lo que ocuparía guardar los IDs como enteros de 32 bits, aunque las listas muy pequeñas (1--16 IDs) pagan hasta 8,5 veces más por metadatos. Las doce consultas evaluadas fueron entre 5 y más de 10\,000 veces más rápidas que recorrer todas las filas, y su coste quedó dominado por el tamaño del resultado.
\end{abstract}
\tableofcontents
\newpage

\section{Problema, pregunta y objetivos}
Buscar repetidamente filas comerciales por palabras de su descripción o por país requiere evaluar predicados sobre muchos registros. Un índice invertido guarda, para cada clave, el conjunto de IDs de las filas que la contienen; así, una consulta AND se resuelve como intersección de conjuntos y una OR como unión. El problema estructural es que las frecuencias son muy heterogéneas: un bitmap completo por clave desperdicia espacio en listas escasas, mientras que un arreglo de IDs no aprovecha las zonas densas.

La pregunta del proyecto es cómo varían el espacio y el tiempo de combinar filtros cuando cambian la cardinalidad de las listas y la densidad dentro de cada bloque de IDs. El objetivo general fue estudiar e implementar el mecanismo de Roaring y evaluarlo en un caso real reproducible. Los objetivos específicos fueron: construir el núcleo y el índice; comprobar los resultados contra un recorrido independiente; medir construcción, espacio y consultas por separado; y explicar ventajas y sobrecostos según cardinalidad y densidad.

El alcance comprende IDs de 32 bits, contenedores array y bitmap, inserción, pertenencia, AND/OR entre todos los pares de contenedores y normalización de resultados. Quedan fuera los run containers, SIMD, concurrencia, eliminación, rank/select y persistencia en disco.

\section{Fundamento y trabajo relacionado}
Los bitmaps clásicos representan un conjunto $S\subseteq\{0,\ldots,N-1\}$ con $N$ bits: las operaciones de conjuntos se vuelven operaciones lógicas palabra a palabra, muy rápidas, pero el espacio es $N/8$ bytes sin importar cuántos elementos tenga $S$. Los esquemas comprimidos por corridas (WAH, Concise) reducen el espacio, pero pierden el acceso aleatorio y obligan a descomprimir secuencialmente. Roaring~\cite{chambi2016} propone una partición en dos niveles: divide el universo de 32 bits en $2^{16}$ bloques de $2^{16}$ valores y, en cada bloque, elige la representación más compacta según cuántos valores contiene. Esto conserva acceso aleatorio, operaciones rápidas en zonas densas y espacio proporcional a la cardinalidad en zonas escasas.

El artículo original incluye además búsquedas adaptativas (galloping) para intersecciones muy desbalanceadas; la biblioteca CRoaring~\cite{lemire2018} añade run containers e instrucciones SIMD. Este proyecto implementa el núcleo sin esas optimizaciones. Como motivación de aplicación, ClickHouse usa Roaring en su índice de texto completo para listas de más de 12 IDs y otra representación para listas menores~\cite{clickhouse2026}, lo que coincide con el sobrecosto de listas pequeñas que observamos en la sección~\ref{sec:espacio}.

\section{El algoritmo Roaring}\label{sec:algoritmo}
\subsection{Descomposición de los enteros}
Cada entero $x$ de 32 bits se divide en dos mitades de 16 bits:
\[
h(x) = x \gg 16, \qquad l(x) = x \;\&\; \texttt{0xFFFF}.
\]
La parte alta $h$ identifica el \emph{bloque} (o contenedor) al que pertenece $x$; la parte baja $l$ es la posición dentro del bloque. Todos los enteros con la misma parte alta comparten un contenedor, y cada contenedor solo necesita almacenar valores de 16 bits. El bitmap Roaring es, por tanto, un vector de pares $(h, C_h)$ ordenado por $h$, donde $C_h$ guarda el conjunto $\{l(x) : x\in S,\ h(x)=h\}$. Los bloques vacíos no se almacenan.

\textbf{Ejemplo.} Para $S=\{3,\ 70\,000,\ 70\,001,\ 131\,077\}$:
\begin{center}\small
\begin{tabular}{rccc}
\toprule
$x$ & $h(x)$ & $l(x)$ & contenedor \\
\midrule
3 & 0 & 3 & $C_0=[3]$ \\
70\,000 & 1 & 4464 & \multirow{2}{*}{$C_1=[4464,4465]$} \\
70\,001 & 1 & 4465 & \\
131\,077 & 2 & 5 & $C_2=[5]$ \\
\bottomrule
\end{tabular}
\end{center}
El conjunto queda como tres contenedores array con claves $0,1,2$. Para reconstruir un ID basta $x = (h \ll 16)\;|\;l$.

\subsection{Dos tipos de contenedor y el umbral 4096}
\begin{itemize}
\item \textbf{Array}: vector ordenado de valores de 16 bits. Ocupa $2c$ bytes para $c$ valores. La búsqueda es binaria.
\item \textbf{Bitmap}: 1024 palabras de 64 bits ($2^{16}$ bits, 8192 bytes). El valor $l$ está presente si el bit $l \bmod 64$ de la palabra $\lfloor l/64\rfloor$ vale 1.
\end{itemize}
El umbral sale de igualar ambos costes: $2c = 8192 \Rightarrow c = 4096$. Con hasta 4096 valores el array es igual o más pequeño; por encima, el bitmap tiene tamaño fijo y se vuelve más compacto. La implementación respeta esta invariante en todo momento: un contenedor con $c\le4096$ es array y uno con $c>4096$ es bitmap. Cada contenedor guarda además su cardinalidad $c$, de modo que la cardinalidad total del conjunto es la suma de los contadores sin recorrer los datos.

\subsection{Inserción y pertenencia}
La inserción localiza el contenedor por búsqueda binaria sobre las claves altas; si no existe, crea un array vacío en la posición ordenada. Dentro de un array, se busca la posición con \texttt{lower\_bound} y se inserta si el valor no estaba; si el contador supera 4096, el array se convierte en bitmap encendiendo el bit de cada uno de sus valores. Dentro de un bitmap, se enciende el bit con una máscara y se incrementa el contador solo si estaba apagado. Los duplicados se ignoran en ambos casos.

La pertenencia (\texttt{contains}) sigue los mismos dos pasos: búsqueda binaria de $h$ entre los $b$ contenedores y luego búsqueda binaria de $l$ en el array o lectura directa del bit en el bitmap.

\subsection{Intersección (AND)}
Las operaciones binarias trabajan en dos niveles. En el primer nivel se recorren los dos vectores de claves ordenadas como en la mezcla de \emph{merge sort}: si las claves difieren se avanza la menor, y si coinciden se combinan los dos contenedores. Para la intersección, una clave presente en un solo operando se descarta, y un contenedor que resulta vacío no se agrega al resultado.

En el segundo nivel, la combinación de dos contenedores usa un método distinto según sus tipos:
\begin{itemize}
\item \textbf{Array AND array}: intersección de dos secuencias ordenadas (\texttt{std::set\_intersection}), lineal en la suma de longitudes. El resultado tiene a lo sumo $\min(c_1,c_2)\le4096$ valores, por lo que siempre es array.
\item \textbf{Array AND bitmap}: se recorre el array y se conserva cada valor cuyo bit esté encendido en el bitmap; cada prueba es $O(1)$. El resultado no puede superar el tamaño del array, así que también es array.
\item \textbf{Bitmap AND bitmap}: se hace AND de las 1024 palabras y se cuentan los bits del resultado. Si la cardinalidad queda en 4096 o menos, el contenedor se \emph{normaliza} a array extrayendo las posiciones de los bits encendidos.
\end{itemize}

\subsection{Unión (OR)}
La unión recorre las claves de la misma forma, pero copia al resultado los contenedores cuya clave aparece en un solo operando. Para claves comunes:
\begin{itemize}
\item \textbf{Array OR array}: unión de secuencias ordenadas (\texttt{std::set\_union}). Si el resultado supera 4096 valores se convierte a bitmap.
\item \textbf{Array OR bitmap}: se copia el bitmap y se enciende el bit de cada valor del array. El resultado es bitmap porque su cardinalidad es al menos la del bitmap ($>4096$).
\item \textbf{Bitmap OR bitmap}: OR de las 1024 palabras y conteo de bits; el resultado siempre es bitmap.
\end{itemize}
Ningún operando se modifica: AND y OR devuelven un bitmap nuevo. Gracias a la normalización, el resultado de cualquier operación cumple la misma invariante de representación que los conjuntos construidos por inserción, por lo que puede volver a combinarse sin casos especiales.

\subsection{Enumeración y conteo de bits}
Para devolver los IDs de una consulta (\texttt{to\_vector}) se recorren los contenedores en orden de clave. Un array se copia directamente; en un bitmap, para cada palabra $w\neq0$ se obtiene la posición del bit menos significativo (conteo de ceros finales), se emite $64i+\text{pos}$ y se apaga ese bit con $w \gets w \;\&\; (w-1)$. El mismo truco ($w\;\&\;(w-1)$ elimina el último 1) se usa para contar bits en AND/OR, de modo que el trabajo es proporcional a los bits encendidos. Como las claves y los valores bajos están ordenados, la salida queda ordenada sin un paso adicional.

\subsection{Complejidad}
\begin{table}[htbp]
\centering\small
\begin{tabular}{ll}
\toprule
Operación & Coste \\
\midrule
Pertenencia (array / bitmap) & $O(\log b + \log c)$ / $O(\log b + 1)$ \\
AND/OR array--array & $O(c_1+c_2)$ \\
AND/OR bitmap--bitmap & $O(1024)$ palabras, sin importar $c$ \\
\bottomrule
\end{tabular}
\caption{Costes de las operaciones principales; $c$ es la cardinalidad de un contenedor y $b$ el número de contenedores.}\label{tab:complejidad}
\end{table}

La consecuencia práctica es que el coste de una operación depende de \emph{cuántos bloques} toca y de \emph{cómo está distribuido} cada bloque, no solo del número total de elementos.

\section{Implementación y aplicación}
\subsection{Núcleo}
La clase \texttt{roaring\_bitmap} (\ruta{include/roaring_bitmap.hpp}, \ruta{src/roaring_bitmap.cpp}) guarda un \texttt{std::vector} de contenedores ordenado por clave. Cada contenedor tiene su clave de 16 bits, su cardinalidad y dos vectores: \texttt{values} (array de \texttt{uint16\_t}) y \texttt{words} (1024 \texttt{uint64\_t}); exactamente uno de ellos está en uso y el tipo se deduce de si \texttt{words} está vacío. La interfaz pública es:
\begin{itemize}
\item \texttt{add(x)} y \texttt{contains(x)}: inserción y pertenencia.
\item \texttt{intersection(o)} y \texttt{union\_with(o)}: AND/OR, devuelven un bitmap nuevo.
\item \texttt{cardinality()}: suma de los contadores de cada contenedor.
\item \texttt{to\_vector()}: enumeración ordenada de los IDs.
\item \texttt{statistics()} y \texttt{memory\_usage()}: tipo, cardinalidad y bytes por contenedor, usados en la evaluación de espacio.
\end{itemize}
Se implementó desde cero, sin usar código de CRoaring ni de otras bibliotecas; los artículos citados son la base conceptual.

\subsection{Índice invertido}
\ruta{src/retail_index.cpp} lee el CSV, normaliza cada fila y construye dos mapas: término $\to$ bitmap y país $\to$ bitmap. El ID de cada fila es su posición en el archivo ($0,\ldots,541908$), por lo que IDs consecutivos caen en el mismo bloque de 65\,536 filas: el universo ocupa solo 9 bloques. Las consultas se traducen directamente a operaciones del núcleo:
\begin{center}\small
\begin{tabular}{ll}
\toprule
Consulta & Operación \\
\midrule
término $t$ & posting$(t)$ \\
país $p$ & posting$(p)$ \\
$t$ AND $p$ & posting$(t)$\texttt{.intersection}(posting$(p)$) \\
$t_1$ OR $t_2$ & posting$(t_1)$\texttt{.union\_with}(posting$(t_2)$) \\
\bottomrule
\end{tabular}
\end{center}
Una clave ausente devuelve un conjunto vacío. Como referencia de corrección se implementó \texttt{scan\_rows}, que evalúa el mismo predicado recorriendo todas las filas normalizadas sin usar el índice. La CLI \ruta{apps/retail_cli.cpp} ejecuta una consulta, muestra las filas originales y, con \verb|--verify|, compara todos los IDs con el scan.

\subsection{Pruebas de corrección}
Las pruebas del núcleo comparan cada operación contra \texttt{std::set} y operaciones ordenadas estándar en conjuntos vacíos, duplicados, extremos de 32 bits, fronteras entre bloques, el paso 4096$\to$4097 y los tres pares de contenedores en AND y OR. En la evaluación, cada consulta compara el vector completo de IDs de Roaring con el del scan antes y después de medir; las doce coincidieron.

\section{Datos y normalización}
Online Retail~\cite{uci2015} contiene transacciones de un comercio británico entre diciembre de 2010 y diciembre de 2011. Se conservaron todas las filas en su orden original, incluidas cancelaciones y descripciones vacías. Los términos son tokens \texttt{[a-z0-9]+} en minúsculas, sin repetición dentro de una fila; los países se pasan a minúsculas con espacios normalizados. No se aplicó stemming ni eliminación de stopwords.

El índice resultante tiene 2310 términos, 38 países y 2\,942\,732 pares fila--clave. Los 2348 postings ocupan 16\,580 contenedores (16\,347 de términos y 233 de países).

\section{Metodología experimental}
Se compiló en modo Release (\texttt{-O3}) y se midió con \texttt{std::chrono::steady\_clock}. La construcción del índice se midió tres veces, aparte de la lectura del CSV y la normalización. Se eligieron doce consultas fijas que cubren casos extremos: términos de frecuencia mínima, mediana y máxima (\texttt{100}: 1 fila, \texttt{peace}: 171, \texttt{set}: 63\,431), el país menos y más frecuente (Saudi Arabia: 10, United Kingdom: 495\,478), tres AND con UK, tres OR entre términos y un término ausente. Cada consulta se midió 15 veces con Roaring y 15 con scan, alternando el orden, y se reporta la mediana. El tiempo de Roaring incluye la materialización del vector de IDs.

El espacio se calcula con un modelo explícito: \emph{payload} ($2c$ bytes por array, 8192 por bitmap), más 2 bytes de clave por contenedor y los metadatos de los objetos (\texttt{sizeof}). Se compara con dos referencias ideales: guardar los IDs como enteros de 32 bits ($4|S|$ bytes) y un bitmap completo por lista ($\lceil541909/8\rceil=67\,739$ bytes). Además se ejecutaron controles sintéticos con semilla fija para aislar el umbral, la dispersión entre bloques y cada par de contenedores.

\section{Resultados}
\subsection{Construcción}
La construcción del índice completo tardó 437 ms (mediana de tres ejecuciones), un coste que se paga una sola vez y que luego se amortiza entre todas las consultas.

\subsection{Espacio}\label{sec:espacio}
\begin{table}[htbp]
\centering\small
\begin{tabular}{lrrrrr}
\toprule
Claves & Payload & Usado & Capacidad & IDs32 & Bitmap/lista \\
\midrule
term & 4\,613\,058 & 5\,583\,930 & 7\,947\,058 & 9\,603\,292 & 156\,477\,090 \\
country & 166\,590 & 180\,550 & 223\,898 & 2\,167\,636 & 2\,574\,082 \\
\bottomrule
\end{tabular}

\caption{Bytes de los postings por tipo de clave. Usado = payload + claves + metadatos; Capacidad incluye las reservas de los vectores.}\label{tab:espacio}
\end{table}

En total los postings usan 5\,764\,480 bytes, el \textbf{49\%} de los 11\,770\,928 bytes que ocuparían como enteros de 32 bits, y 28 veces menos que un bitmap completo por lista (159 MB). El ahorro viene de dos efectos del algoritmo: cada ID se guarda con 16 bits en lugar de 32 dentro de los arrays, y los bloques densos (por ejemplo, UK) pasan a bitmaps de tamaño fijo.

\begin{figure}[htbp]
\centering
\begin{tikzpicture}
\begin{axis}[ybar,bar width=12pt,width=0.9\textwidth,height=6cm,ymin=0,ymax=11.5,enlarge x limits=0.15,
ylabel={Bytes Roaring / bytes IDs32},xlabel={Cardinalidad del posting},
symbolic x coords={a,b,c,d},xtick=data,
xticklabels={1--16,17--256,257--4096,$>4096$},
nodes near coords,nodes near coords style={font=\scriptsize,/pgf/number format/fixed,/pgf/number format/precision=2},
legend style={at={(0.98,0.97)},anchor=north east,font=\small},
ymajorgrids,grid style={gray!20},axis lines*=left,
extra y ticks={1},extra y tick style={grid=major,major grid style={black,dashed}}]
\addplot[fill=serieroaring,draw=none] coordinates {(a,8.496759822) (b,1.519760317) (c,0.622734497) (d,0.479418480)};
\addplot[fill=serieroaring!35,draw=none] coordinates {(a,9.953422438) (b,2.242846005) (c,0.924031758) (d,0.658488922)};
\legend{Usado,Capacidad}
\end{axis}
\end{tikzpicture}
\caption{Espacio de los postings de términos relativo a guardar los IDs como enteros de 32 bits, por grupo de cardinalidad. La línea punteada marca la igualdad: por debajo, Roaring ahorra; por encima, cuesta más.}\label{fig:espacio}
\end{figure}

\begin{table}[htbp]
\centering\small
\begin{tabular}{lrrrrr}
\toprule
Cardinalidad & Listas & Usado (B) & Capacidad (B) & IDs32 (B) & U/IDs32 \\
\midrule
1-16 & 465 & 83\,914 & 98\,300 & 9\,876 & 8.50 \\
17-256 & 909 & 608\,202 & 897\,578 & 400\,196 & 1.52 \\
257-4096 & 805 & 2\,104\,870 & 3\,123\,268 & 3\,380\,044 & 0.62 \\
$>$4096 & 131 & 2\,786\,944 & 3\,827\,912 & 5\,813\,176 & 0.48 \\
\bottomrule
\end{tabular}

\caption{Postings de términos agrupados por cardinalidad.}\label{tab:grupos}
\end{table}

Ese ahorro agregado esconde un sobrecosto en listas pequeñas (figura~\ref{fig:espacio}, tabla~\ref{tab:grupos}). Las 465 listas de 1--16 IDs ocupan 8,5 veces más que sus IDs de 32 bits: con tan pocos elementos, el payload es mínimo y dominan la clave y los metadatos del contenedor (dos vectores y un contador). El grupo 17--256 también cuesta más. A partir de 257 IDs Roaring ahorra, y el ahorro crece con la cardinalidad.

\begin{figure}[htbp]
\centering
\begin{tikzpicture}
\begin{axis}[ybar,bar width=12pt,width=0.85\textwidth,height=6cm,ymode=log,ymin=1,ymax=60000,
log origin=infty,
ylabel={Contenedores (escala log)},xlabel={Densidad del bloque ($c/65536$)},
symbolic x coords={a,b,c,d},xtick={a,b,c,d},
xticklabels={$\le1\%$,1--6{,}25\%,6{,}25--25\%,$>25\%$},
point meta=rawy,nodes near coords,nodes near coords style={font=\scriptsize},
enlarge x limits=0.18,
legend style={at={(0.98,0.97)},anchor=north east,font=\small},
ymajorgrids,grid style={gray!20},axis lines*=left]
\addplot[fill=serieroaring,draw=none] coordinates {(a,15558) (b,733) (c,56)};
\addplot[fill=seriescan,draw=none] coordinates {(a,201) (b,23) (d,9)};
\legend{Términos,Países}
\end{axis}
\end{tikzpicture}
\caption{Distribución de los contenedores según su densidad. Solo los bloques con más de 4096 valores ($>6{,}25\%$) son bitmaps: 56 de términos y 9 de países.}\label{fig:densidad}
\end{figure}

La figura~\ref{fig:densidad} muestra por qué casi todo el índice está formado por arrays: el 95\% de los bloques de términos tiene como máximo un 1\% de densidad. Solo 65 de los 16\,580 contenedores son bitmaps, pero concentran gran parte de los IDs: los 9 bitmaps de UK guardan 495\,478 IDs en 73\,728 bytes (1,2 bits por ID). Es la densidad local del bloque, y no el tamaño total de la lista, lo que determina la representación.

\subsection{Consultas}
\begin{table}[htbp]
\centering\small
\begin{tabular}{lrrrr}
\toprule
Consulta & IDs & Scan (ms) & Roaring (ms) & Veces más rápido \\
\midrule
100 & 1 & 16.527 & 0.0029 & 5\,791 \\
peace & 171 & 20.615 & 0.0012 & 16\,843 \\
set & 63\,431 & 21.745 & 0.1974 & 110 \\
Saudi Arabia & 10 & 2.937 & 0.0003 & 10\,839 \\
United Kingdom & 495\,478 & 4.390 & 0.8721 & 5 \\
100 AND UK & 1 & 17.769 & 0.0007 & 26\,601 \\
peace AND UK & 155 & 20.744 & 0.0091 & 2\,291 \\
set AND UK & 56\,118 & 18.422 & 0.2371 & 78 \\
100 OR peace & 172 & 27.176 & 0.0073 & 3\,725 \\
100 OR set & 63\,432 & 26.983 & 0.1796 & 150 \\
peace OR set & 63\,602 & 29.089 & 0.1844 & 158 \\
ausente & 0 & 11.937 & 0.0002 & 48\,134 \\
\bottomrule
\end{tabular}
\caption{Tiempo mediano de 15 repeticiones por consulta. UK = United Kingdom. ``Veces más rápido'' = tiempo del scan dividido entre tiempo de Roaring.}\label{tab:latencias}
\end{table}

\begin{figure}[htbp]
\centering
\begin{tikzpicture}
\begin{axis}[xbar,bar width=4pt,width=0.82\textwidth,height=10cm,
xmode=log,xmin=0.0001,xmax=100,log origin=infty,
xlabel={Mediana (ms, escala log)},
symbolic y coords={ausente,peace OR set,100 OR set,100 OR peace,set AND UK,peace AND UK,100 AND UK,United Kingdom,Saudi Arabia,set,peace,100},
ytick=data,yticklabel style={font=\small},
enlarge y limits=0.05,
legend style={at={(0.5,1.02)},anchor=south,legend columns=2,font=\small},
xmajorgrids,grid style={gray!20},axis lines*=left]
\addplot[fill=seriescan,draw=none] coordinates {(16.526909,100) (20.615261,peace) (21.744897,set) (2.937224,Saudi Arabia) (4.389661,United Kingdom) (17.769275,100 AND UK) (20.744259,peace AND UK) (18.422302,set AND UK) (27.176090,100 OR peace) (26.983176,100 OR set) (29.088881,peace OR set) (11.937236,ausente)};
\addplot[fill=serieroaring,draw=none] coordinates {(0.002854,100) (0.001224,peace) (0.197382,set) (0.000271,Saudi Arabia) (0.872068,United Kingdom) (0.000668,100 AND UK) (0.009054,peace AND UK) (0.237067,set AND UK) (0.007296,100 OR peace) (0.179648,100 OR set) (0.184386,peace OR set) (0.000248,ausente)};
\legend{Scan,Roaring}
\end{axis}
\end{tikzpicture}
\caption{Latencia mediana de cada consulta con el índice Roaring y con el recorrido completo de filas (scan). El eje es logarítmico: cada división es un factor 10.}\label{fig:latencias}
\end{figure}

Las doce consultas fueron más rápidas con el índice (figura~\ref{fig:latencias}). El scan cuesta siempre entre 3 y 29 ms porque recorre las 541\,909 filas sin importar el resultado. Roaring, en cambio, solo toca los contenedores de las claves involucradas, y su tiempo va de 0,25 $\mu$s (término ausente) a 0,87 ms (UK).

\begin{figure}[htbp]
\centering
\begin{tikzpicture}
\begin{axis}[width=0.8\textwidth,height=6.5cm,xmode=log,ymode=log,
xmin=0.5,xmax=1e6,ymin=1,ymax=1e5,
xlabel={Filas devueltas por la consulta (escala log)},ylabel={Veces más rápido que el scan},
xmajorgrids,ymajorgrids,grid style={gray!20},axis lines*=left]
\addplot[only marks,mark=*,mark size=3pt,color=serieroaring,
nodes near coords,point meta=explicit symbolic,
nodes near coords style={font=\scriptsize,color=black,anchor=west,xshift=2pt}]
coordinates {(1,5790.8)[100] (171,16842.5)[peace] (63431,110.2)[] (10,10838.5)[Saudi Arabia] (495478,5.0)[UK] (1,26600.7)[100 AND UK] (155,2291.2)[] (56118,77.7)[] (172,3724.8)[] (63432,150.2)[] (63602,157.8)[set, set AND UK, OR con set]};
\end{axis}
\end{tikzpicture}
\caption{Cada punto es una consulta. Cuanto más a la derecha, más filas devuelve; cuanto más arriba, mayor es la ventaja de Roaring sobre el scan.}\label{fig:aceleracion}
\end{figure}

La figura~\ref{fig:aceleracion} responde a una pregunta sencilla: \emph{¿de qué depende cuánto gana el índice?} Cada punto es una de las consultas; en el eje horizontal está cuántas filas devuelve y en el vertical cuántas veces más rápida fue con Roaring que con el scan. Los puntos bajan de izquierda a derecha: las consultas que devuelven pocas filas (\texttt{100}, \texttt{peace}, Saudi Arabia) son miles de veces más rápidas, las que devuelven decenas de miles (\texttt{set} y sus combinaciones) son unas 100 veces más rápidas, y UK, que devuelve casi todas las filas, solo 5 veces.

La razón es que el scan siempre hace el mismo trabajo (revisar las 541\,909 filas), mientras que Roaring hace un trabajo proporcional a lo que devuelve: con pocas filas solo toca uno o dos contenedores pequeños, pero con muchas tiene que recorrer los contenedores y copiar cada ID al vector de salida. Por eso la conclusión principal sobre tiempo es que \textbf{el coste de Roaring depende sobre todo del tamaño del resultado}, no de la operación. Por ejemplo, \texttt{set AND UK} (56\,118 filas) tarda 0,24 ms, casi lo mismo que devolver solo \texttt{set} (0,20 ms): la intersección en sí es barata. Los tiempos por debajo de 1 $\mu$s están cerca de la resolución del reloj, así que las ventajas más extremas deben leerse como órdenes de magnitud, no como valores exactos.

\subsection{Controles sintéticos}
\begin{table}[htbp]
\centering\scriptsize
\begin{tabular}{llrrrrr}
\toprule
Caso & Operacion & IDs & A & B & Payload (B) & Mediana ($\mu$s) \\
\midrule
AA & and & 1\,000 & 1 & 0 & 2\,000 & 2.399 \\
AA & or & 5\,000 & 0 & 1 & 8\,192 & 13.779 \\
BB vacio & and & 0 & 0 & 0 & 0 & 2.623 \\
BB a array & and & 1\,000 & 1 & 0 & 2\,000 & 3.731 \\
8192 bloques & build & 8\,192 & 8192 & 0 & 16\,384 & 48258.960 \\
bloque 4096 & build & 4\,096 & 1 & 0 & 8\,192 & 282.052 \\
bloque 4097 & build & 4\,097 & 0 & 1 & 8\,192 & 284.974 \\
bloque 8192 & build & 8\,192 & 0 & 1 & 8\,192 & 303.267 \\
4096 a 4097 & insert & 4\,097 & 0 & 1 & 8\,192 & 4.779 \\
\bottomrule
\end{tabular}

\caption{Selección de escenarios sintéticos; A/B = contenedores array/bitmap del resultado. Aquí el tiempo no incluye la materialización de IDs.}\label{tab:sinteticos}
\end{table}

Los sintéticos (tabla~\ref{tab:sinteticos}) confirman el comportamiento descrito en la sección~\ref{sec:algoritmo}:
\begin{itemize}
\item \textbf{Umbral.} Un bloque con 4096 valores (array) y otro con 4097 (bitmap) ocupan ambos 8192 bytes: el umbral cambia la representación sin un salto de espacio. La inserción que provoca la conversión cuesta 4,8 $\mu$s, un coste puntual.
\item \textbf{Normalización.} La unión de dos arrays puede producir un bitmap (5000 IDs) y la intersección de dos bitmaps puede volver a array (1000 IDs) o quedar vacía, en cuyo caso el contenedor se elimina.
\item \textbf{Bitmap--bitmap.} AND y OR entre bitmaps tardan 3--4 $\mu$s independientemente de la cardinalidad, porque siempre recorren 1024 palabras.
\item \textbf{Distribución.} Los mismos 8192 IDs concentrados en un bloque se construyen en 0,30 ms y ocupan 8192 bytes; dispersos en 8192 bloques (uno por bloque) ocupan el doble de payload, crean 8192 contenedores y tardan 48 ms, unas 160 veces más. El coste adicional proviene de crear un contenedor por ID y de insertar cada clave nueva en su posición ordenada dentro del vector de contenedores.
\end{itemize}

\subsection{Recuperación de filas}
La consulta \texttt{heart AND United Kingdom} con la CLI devuelve 37\,513 IDs idénticos a los del scan:
\begin{verbatim}
./build/release/retail_cli \
  --data data/prepared/online_retail.csv \
  --type and --first heart --second 'United Kingdom' \
  --limit 5 --verify
\end{verbatim}
Los primeros IDs son 0, 4, 19, 49 y 62; por ejemplo, la fila 0 corresponde a la factura 536365, \texttt{WHITE HANGING HEART T-LIGHT HOLDER}.

\section{Discusión y límites}
\textbf{Qué muestra la evidencia.} La combinación array/bitmap explica tanto el ahorro agregado de espacio como el sobrecosto de las listas pequeñas, donde los metadatos dominan. Las operaciones entre contenedores se comportan según su análisis: bitmap--bitmap tiene coste constante, array--array es lineal y la normalización mantiene la invariante en los resultados. Con el índice construido, los filtros repetidos son mucho más rápidos que recorrer las filas, y el coste restante depende sobre todo de cuántos IDs se devuelven.

\textbf{Qué no muestra.} El scan no es un índice alternativo, así que el experimento separa el beneficio de \emph{tener} un índice, no el de elegir Roaring frente a listas ordenadas o bitmaps completos; esas alternativas solo se comparan en espacio. Doce consultas en un solo dataset y un solo equipo no representan una carga real, y el modelo de memoria no incluye el allocator ni el diccionario de claves.

\section{Mejoras futuras}
Ninguna de estas mejoras está implementada:
\begin{itemize}
\item \textbf{Comparar representaciones con el mismo contrato}: listas ordenadas de IDs32 y bitmaps completos, midiendo tiempo de consulta, para aislar el efecto de Roaring frente al de indexar.
\item \textbf{Construcción en lote}: como los IDs llegan en orden creciente, se puede agregar al final del contenedor en vez de insertar con desplazamiento y reservar el tamaño final; la diferencia entre capacidad y espacio usado es de 2,4 MB.
\item \textbf{Listas pequeñas inline}: guardar las listas de pocos IDs directamente como enteros, evitando los metadatos de contenedor, con un umbral calibrado para este índice.
\item \textbf{Run containers e intersección adaptativa}: compresión de corridas de IDs consecutivos y búsqueda galloping para AND entre arrays muy desbalanceados~\cite{chambi2016,lemire2018}.
\end{itemize}

\section{Conclusiones}
Se implementó desde cero el núcleo de Roaring Bitmaps (descomposición en clave alta y valor bajo, contenedores array y bitmap con umbral 4096, inserción con conversión, pertenencia, AND/OR para los tres pares de contenedores y normalización) y se aplicó a un índice invertido sobre todas las filas de Online Retail, con resultados idénticos a un recorrido completo.

En espacio, los postings ocupan la mitad que los IDs de 32 bits, pero las listas pequeñas cuestan hasta 8,5 veces más por los metadatos. En tiempo, el índice responde entre 5 y más de 10\,000 veces más rápido que el scan, y el coste está dominado por el tamaño del resultado. Los controles sintéticos muestran que la densidad dentro de cada bloque y la distribución de los IDs entre bloques importan más que la cardinalidad total del conjunto.

\appendix
\section{Reproducción}
Desde la raíz del repositorio:
\begin{verbatim}
python3 scripts/prepare_data.py
cmake -S . -B build/release -DCMAKE_BUILD_TYPE=Release
cmake --build build/release -j2
ctest --test-dir build/release --output-on-failure
python3 scripts/run_experiments.py \
  --binary build/release/retail_benchmark \
  --data data/prepared/online_retail.csv \
  --output results/nueva_ejecucion
python3 scripts/summarize_results.py results/nueva_ejecucion
\end{verbatim}
El directorio de salida debe ser nuevo. Si falta el XLSX, \verb|prepare_data.py --download| obtiene el archivo oficial.

\begin{thebibliography}{9}
\bibitem{chambi2016}
S. Chambi, D. Lemire, O. Kaser y R. Godin. ``Better bitmap performance with Roaring bitmaps.'' \emph{Software: Practice and Experience}, 46(5), 709--719, 2016. \href{https://doi.org/10.1002/spe.2325}{doi:10.1002/spe.2325}.
\bibitem{lemire2018}
D. Lemire, O. Kaser, N. Kurz, L. Deri, C. O'Hara, F. Saint-Jacques y G. Ssi-Yan-Kai. ``Roaring bitmaps: Implementation of an optimized software library.'' \emph{Software: Practice and Experience}, 48(4), 867--895, 2018. \href{https://doi.org/10.1002/spe.2560}{doi:10.1002/spe.2560}.
\bibitem{uci2015}
D. Chen. \emph{Online Retail} [conjunto de datos]. UCI Machine Learning Repository, 2015. \href{https://doi.org/10.24432/C5BW33}{doi:10.24432/C5BW33}. CC BY 4.0.
\bibitem{clickhouse2026}
E. Ahmadov, J. Aguilar y G. Larionov. ``Building high-performance full-text search for object storage.'' ClickHouse Engineering, 2026. \url{https://clickhouse.com/blog/clickhouse-full-text-search-object-storage}.
\end{thebibliography}
\end{document}
